\chapter{Support Vector Machines}

\section{Margin Maximization}

In previous chapters, we studied linear classifiers and their extensions via feature maps. Support Vector Machines (SVMs) refine these ideas by introducing a geometric principle: \emph{margin maximization}.

\medskip

\noindent
\textbf{Binary classification setting.}  
Let $y_i \in \{-1, +1\}$ and consider a linear classifier
\[
f(x) = \langle w, x \rangle + b.
\]

\medskip

\noindent
A point is correctly classified if
\[
y_i (\langle w, x_i \rangle + b) > 0.
\]

\medskip

\noindent
\textbf{Margin.}  
The (functional) margin of a data point is defined as
\[
y_i (\langle w, x_i \rangle + b).
\]

\medskip

\noindent
\textbf{Geometric margin.}  
The distance of a point to the decision boundary is
\[
\frac{y_i (\langle w, x_i \rangle + b)}{\|w\|}.
\]

\medskip

\noindent
\textbf{Goal.}  
Find a classifier that maximizes the minimum margin over all training points.

\medskip

\noindent
\textbf{Optimization problem.}
\[
\max_{w,b} \; \min_{i} \; \frac{y_i (\langle w, x_i \rangle + b)}{\|w\|}.
\]

\medskip

\noindent
\textbf{Equivalent formulation.}  
By fixing the scale, we obtain
\[
\min_{w,b} \; \frac{1}{2}\|w\|^2
\quad \text{subject to} \quad
y_i (\langle w, x_i \rangle + b) \geq 1 \quad \forall i.
\]

\medskip

\noindent
\textbf{Interpretation.}
\begin{itemize}
    \item Maximizing margin $\leftrightarrow$ minimizing $\|w\|$
    \item Encourages robust separation between classes
\end{itemize}

\medskip

\noindent
\textbf{Support vectors.}  
Only points that lie closest to the decision boundary (i.e., those with equality in the constraint) influence the solution.

\medskip

\noindent
\textbf{Insight.}  
SVMs focus on the most informative points, rather than all training data.

\section{Hard and Soft Margins}

The formulation above assumes that the data is linearly separable. In practice, this is rarely the case.

\medskip

\noindent
\textbf{Hard-margin SVM.}
\begin{itemize}
    \item Assumes perfect separability
    \item Enforces constraints:
    \[
    y_i (\langle w, x_i \rangle + b) \geq 1
    \]
\end{itemize}

\medskip

\noindent
\textbf{Limitation.}  
Sensitive to noise and outliers; may have no feasible solution.

\medskip

\noindent
\textbf{Soft-margin SVM.}  
Introduce slack variables $\xi_i \geq 0$:
\[
y_i (\langle w, x_i \rangle + b) \geq 1 - \xi_i.
\]

\medskip

\noindent
\textbf{Objective.}
\[
\min_{w,b,\xi} \; \frac{1}{2}\|w\|^2 + C \sum_{i=1}^n \xi_i,
\]
where $C > 0$ controls the tradeoff between margin size and classification error.

\medskip

\noindent
\textbf{Interpretation.}
\begin{itemize}
    \item $\xi_i$ measures violation of the margin constraint
    \item Larger $C$ $\rightarrow$ penalizes errors more strongly
    \item Smaller $C$ $\rightarrow$ allows more flexibility
\end{itemize}

\medskip

\noindent
\textbf{Connection to loss functions.}  
The soft-margin formulation corresponds to minimizing a \emph{hinge loss}:
\[
\ell(x_i, y_i) = \max\big(0, 1 - y_i f(x_i)\big).
\]

\medskip

\noindent
\textbf{Insight.}  
SVMs combine margin maximization with robustness to non-separable data.

\section{Dual Formulation (Conceptual)}

The SVM optimization problem admits an alternative representation known as the \emph{dual formulation}.

\medskip

\noindent
\textbf{Key idea.}  
Instead of optimizing over $w$, we introduce variables $\alpha_i$ associated with each data point.

\medskip

\noindent
\textbf{Dual problem (informal).}
\[
\max_{\alpha} \; \sum_{i=1}^n \alpha_i - \frac{1}{2} \sum_{i,j} \alpha_i \alpha_j y_i y_j \langle x_i, x_j \rangle,
\]
subject to constraints on $\alpha_i$.

\medskip

\noindent
\textbf{Important observations.}
\begin{itemize}
    \item The solution depends only on inner products $\langle x_i, x_j \rangle$
    \item Only support vectors have nonzero $\alpha_i$
\end{itemize}

\medskip

\noindent
\textbf{Reconstruction of $w$.}
\[
w = \sum_{i=1}^n \alpha_i y_i x_i.
\]

\medskip

\noindent
\textbf{Insight.}  
The classifier is expressed as a weighted combination of training points.

\medskip

\noindent
\textbf{Practical significance.}
\begin{itemize}
    \item Enables the use of kernel methods
    \item Reveals sparsity in the solution
\end{itemize}

\section{Kernel Methods Revisited}

The dual formulation allows us to extend SVMs to nonlinear classification.

\medskip

\noindent
\textbf{Kernel substitution.}  
Replace inner products with a kernel function:
\[
\langle x_i, x_j \rangle \;\rightarrow\; k(x_i, x_j).
\]

\medskip

\noindent
\textbf{Resulting classifier.}
\[
f(x) = \sum_{i=1}^n \alpha_i y_i k(x_i, x) + b.
\]

\medskip

\noindent
\textbf{Interpretation.}
\begin{itemize}
    \item Implicitly maps data into a feature space
    \item Constructs nonlinear decision boundaries in the original space
\end{itemize}

\medskip

\noindent
\textbf{Examples.}
\begin{itemize}
    \item Polynomial kernel
    \item Gaussian (RBF) kernel
\end{itemize}

\medskip

\noindent
\textbf{Geometric view.}  
Data is transformed into a space where it becomes linearly separable, and a maximum-margin separator is found there.

\medskip

\noindent
\textbf{Advantages.}
\begin{itemize}
    \item Flexible nonlinear modeling
    \item Strong theoretical guarantees
\end{itemize}

\medskip

\noindent
\textbf{Limitations.}
\begin{itemize}
    \item Computational cost grows with dataset size
    \item Less scalable than modern deep learning approaches
\end{itemize}

\section{A Unifying Perspective}

Support Vector Machines integrate several core ideas:

\begin{itemize}
    \item \textbf{Geometry:} margin maximization
    \item \textbf{Optimization:} convex quadratic programming
    \item \textbf{Representation:} feature maps and kernels
\end{itemize}

\medskip

\noindent
\textbf{Key insights.}
\begin{itemize}
    \item Good classifiers maximize separation between classes
    \item Only a subset of data points determines the solution
    \item Nonlinearity can be introduced through kernels
\end{itemize}

\medskip

\noindent
SVMs illustrate how combining geometry and optimization leads to powerful learning algorithms.

\section{Outlook}

This chapter introduced support vector machines as a principled approach to classification:
\begin{itemize}
    \item margin-based learning,
    \item soft constraints for real-world data,
    \item dual formulation and sparsity,
    \item kernel methods for nonlinear models.
\end{itemize}

\medskip

In the next chapter, we return to model evaluation and selection, focusing on how to compare models and choose hyperparameters in practice.

\medskip

\noindent
\textbf{Guiding principle.}  
Effective learning balances expressivity with robustness, often through geometric and optimization principles.